\documentclass[10pt]{article}

% =========================================================
% Packages
% =========================================================
\usepackage{xltxtra}
\usepackage{bookmark}
\usepackage{hyperref}
\hypersetup{hidelinks}

\usepackage{url}
\urlstyle{tt}

\usepackage{multicol}
\usepackage{xcolor}
\usepackage{calc}
\usepackage{graphicx}
\usepackage{tikz}
\usepackage{eso-pic}
\usetikzlibrary{calc}

\usepackage{fontspec}
\usepackage{xeCJK}
\usepackage{relsize}
\usepackage{xspace}
\usepackage{fontawesome}
\usepackage{titlesec}
\usepackage{enumitem}
\usepackage{siunitx}
\usepackage{amssymb}
\usepackage{tabularx}
\usepackage{array}

% =========================================================
% 页面设置
% =========================================================
\usepackage[
    a4paper,
    left=0.65cm,
    right=0.65cm,
    top=1.78cm,
    bottom=1.28cm,
    nohead
]{geometry}

% =========================================================
% 全局排版
% =========================================================

\setlength{\parindent}{0pt}
\setlength{\parskip}{0pt}

\pagenumbering{gobble}

% =========================================================
% 中文排版
% =========================================================

% 中文与英文/数字之间不额外增加距离
\CJKsetecglue{}

% 不人为增加中文字距
\renewcommand{\CJKglue}{\hskip 0pt}

% 防止 TeX 强行拉开中文字距
\emergencystretch=0pt

% =========================================================
% 字体
% =========================================================

\setmainfont[
    Path=fonts/,
    Extension=.ttf,
    BoldFont=* Bold,
]{Microsoft Yahei}

\setCJKmainfont[
    Path=fonts/,
    Extension=.ttf,
    BoldFont=* Bold,
]{Microsoft Yahei}

% =========================================================
% 行距 / 表格 / 多栏
% =========================================================

% 正文行距
\linespread{1.06}

% 表格行距
\renewcommand{\arraystretch}{1.0}

% 表格列间距
\setlength{\tabcolsep}{3pt}

% 多栏上下间距
\setlength{\multicolsep}{0.05em}

% 两栏之间的间距
\setlength{\columnsep}{1.5em}

% =========================================================
% 列表
% =========================================================

\setlist[itemize]{
    topsep=0.08em,
    partopsep=0pt,
    leftmargin=1.4em,
    itemsep=0.05em,
    parsep=0pt
}

\setlist[enumerate]{
    topsep=0.08em,
    partopsep=0pt,
    leftmargin=1.4em,
    itemsep=0.05em,
    parsep=0pt
}

% =========================================================
% 主题颜色
% =========================================================

\definecolor{CHU_Blue}{RGB}{19,57,132}

% =========================================================
% 一级标题
% =========================================================

\titleformat{\section}
    {\fontsize{13.2pt}{14.2pt}\selectfont
     \bfseries
     \raggedright}
    {}
    {0em}
    {}
    [
        {\color{CHU_Blue}\titlerule}
    ]

\titlespacing*{\section}
    {0cm}
    {0.55em}
    {0.22em}

% =========================================================
% 二级标题
% =========================================================

\titleformat{\subsection}
    {\fontsize{10.3pt}{11.4pt}\selectfont
     \bfseries
     \raggedright}
    {}
    {0em}
    {}
    []

\titlespacing*{\subsection}
    {0cm}
    {0.42em}
    {0.12em}

% =========================================================
% 学院名称
% =========================================================

\newcommand{\school}{长安都柏林国际交通学院 | CDIC}

% =========================================================
% 页脚联系方式
% =========================================================

\newcommand{\contact}{%
    {%
        \scriptsize
        \color{white}%
        \faEnvelope\quad
        \href{mailto:usagi@chiikawa.world}
        {usagi@chiikawa.world}%
        \hspace{4em}%
        \faWechat\quad
        USAGI\_YAHA%
        \hspace{4em}%
        \faPhone\quad
        YAHA-URA-2026%
        \hspace{4em}%
        \faGithub\quad
        \href{https://github.com/zhangkeo1111-gif/}
        {github.com/zhangkeo1111-gif}%
    }%
}

% =========================================================
% 页眉 / 页脚 / 水印
% =========================================================
%
% 保持 GitHub 原模板效果：
% 页眉和页脚图片只指定 width=\paperwidth，
% 不人为改变图片纵横比。
%

\newcommand{\resumePageStyle}{%
    \begin{tikzpicture}[remember picture, overlay]

        % -------------------------------------------------
        % 页眉背景
        % -------------------------------------------------

        \node[
            anchor=north,
            inner sep=0pt
        ] (header) at (current page.north) {%
            \includegraphics[
                width=\paperwidth
            ]{images/Head_chu_final.png}%
        };

        % -------------------------------------------------
        % 左侧学校 Logo
        % -------------------------------------------------

        \node[
            anchor=west
        ] (school_logo) at (header.west) {%
            \hspace{0.5cm}%
            \includegraphics[
                width=0.25\textwidth
            ]{images/chu_logo_white.png}%
        };

        % -------------------------------------------------
        % 右侧学院名称
        %
        % 必须显式固定字号，
        % 防止 shipout 时继承正文 / section 的当前字号。
        % -------------------------------------------------

        \node[
            anchor=east
        ] (school_name) at (header.east) {%
            {%
                \color{white}%
                \fontsize{11pt}{12pt}\selectfont
                \bfseries
                \school
            }%
            \hspace{0.5cm}%
        };

        % -------------------------------------------------
        % 页脚
        % -------------------------------------------------

        \node[
            anchor=south,
            inner sep=0pt
        ] (footer) at (current page.south) {%
            \includegraphics[
                width=\paperwidth
            ]{images/Foot_chu_final}%
        };

        % 页脚联系方式
        \node[
            anchor=center
        ] at (footer.center) {%
            \contact
        };

        % -------------------------------------------------
        % 背景水印
        % -------------------------------------------------

        \node[
            opacity=0.1
        ] at (current page.center) {%
            \includegraphics[
                width=0.45\paperwidth,
                keepaspectratio
            ]{images/chu_logo_big.jpg}%
        };

    \end{tikzpicture}%
}

\AddToShipoutPictureFG{\resumePageStyle}

% =========================================================
% Document
% =========================================================

\begin{document}

% =========================================================
% 个人信息
% =========================================================

\noindent
\begin{minipage}[c]{0.82\textwidth}

    \section{
        \makebox[\widthof{\faUser}][c]{
            \color{CHU_Blue}{\faUser}
        }
        \quad
        个人信息
    }

    \vspace{-0.05em}

    \begin{tabularx}
        {\linewidth}
        {
            @{}
            p{\widthof{出生日期:}}
            X
            p{\widthof{政治面貌:}}
            X
            @{}
        }

        姓\qquad 名:
        &
        乌萨奇
        &
        性\qquad 别:
        &
        不详
        \\

        年\qquad 龄:
        &
        不详
        &
        政治面貌:
        &
        自由快乐小兔
        \\

        联系方式:
        &
        YAHA-URA-2026
        &
        电子邮箱:
        &
        \href{mailto:usagi@chiikawa.world}
        {usagi@chiikawa.world}
        \\

    \end{tabularx}

\end{minipage}
\hfill
\begin{minipage}[c]{0.16\textwidth}

    \centering

    \includegraphics[
        width=0.90\linewidth
    ]{images/Photo_usagi.jpg}

\end{minipage}

\vspace{-0.20em}

% =========================================================
% 教育背景
% =========================================================

\section{
    \makebox[\widthof{\faGraduationCap}][c]{
        \color{CHU_Blue}{\faGraduationCap}
    }
    \quad
    教育背景
}

\noindent
\begin{tabular*}
    {\textwidth}
    {
        @{}
        l
        @{\extracolsep{\fill}}
        l
        @{\extracolsep{\fill}}
        r
        @{}
    }

    吉伊卡哇世界劳动学院
    &
    草丛采集与讨伐工程
    &
    2025年6月-至今
    \\

    \textbf{GPA: 5.0/5.0}
    &
    \textbf{专业排名: 1/999}
    &
    \textbf{综测排名: 1/999}
    \\

\end{tabular*}

\vspace{0em}

\begin{itemize}

    \item
    核心课程：草丛采集基础、讨伐任务导论、洞穴探索安全、
    异常生物识别、零食资源发现、危机快速脱离、
    小队协作、即兴口号表达、强运行为学。

\end{itemize}

% =========================================================
% 所获荣誉
% =========================================================

\section{
    \makebox[\widthof{\faStar}][c]{
        \color{CHU_Blue}{\faStar}
    }
    \quad
    所获荣誉
}

\begin{multicols}{2}

\begin{itemize}

    \item 吉伊卡哇世界高机动行动派代表

    \item 多次完成草丛采集与讨伐劳动任务

    \item “YAHA/URA”式高辨识度沟通达人

    \item 小队突发事件快速反应核心成员

    \item 零食发现与即时快乐转化优秀个体

    \item 洞穴、草丛与街区复合场景适应能手

\end{itemize}

\end{multicols}

% =========================================================
% 竞赛经历
% =========================================================

\section{
    \makebox[\widthof{\faTrophy}][c]{
        \color{CHU_Blue}{\faTrophy}
    }
    \quad
    竞赛经历
}

\noindent
\begin{tabular*}
    {\textwidth}
    {
        @{}
        l
        @{\extracolsep{\fill}}
        l
        l
        r
        @{}
    }

    吉伊卡哇世界草丛采集挑战赛
    &
    突击成员
    &
    国家级一等奖
    &
    2026年6月
    \\

    吉伊卡哇世界讨伐委托实战赛
    &
    高机动输出
    &
    国家级一等奖
    &
    2026年6月
    \\

    三人小队协同探索赛
    &
    核心队员
    &
    国家级一等奖
    &
    2026年6月
    \\

    零食发现与快速消费竞赛
    &
    个人选手
    &
    国家级一等奖
    &
    2026年6月
    \\

\end{tabular*}

% =========================================================
% 实习经历
% =========================================================

\section{
    \makebox[\widthof{\faGears}][c]{
        \color{CHU_Blue}{\faGears}
    }
    \quad
    实习经历
}

\subsection{
    劳动铠甲事务所
    \hfill
    讨伐与采集实习兔
    \hfill
    2025年6月-2026年6月
}

\hspace*{1.5em}
以草丛采集、街区巡游与小型讨伐任务为核心工作内容，
长期活跃于吉伊卡哇世界的一线劳动场景。
实习期间，熟悉了从任务领取、风险识别、目标接近到战利品回收的完整流程，
能够在信息不足、环境不稳定和同伴情绪波动的情况下保持极高行动速度。
面对突发异常事件时，
能够通过“YAHA”“URA”等高频短指令完成快速沟通，
并以跳跃、挥动、冲刺等方式带动小队推进。
该经历使乌萨奇形成了强烈的现场直觉、
极高的危险适应性与无法被常规绩效表格完全解释的执行风格。

% =========================================================
% 项目经历
% =========================================================

\section{
    \makebox[\widthof{\faGears}][c]{
        \color{CHU_Blue}{\faGears}
    }
    \quad
    项目经历
}

% ---------------------------------------------------------
% 项目 1
% ---------------------------------------------------------

\subsection{
    高机动草丛检索与异常目标快速接近算法
    \hfill
    日常劳动项目
    \hfill
    2025年6月-2026年6月
}

\hspace*{1.5em}
作为小队中的高机动行动单元，
围绕草丛环境下的资源发现与异常目标规避问题，
构建了以直觉判断、快速接近和即时撤离为核心的行动策略。
在任务执行中，
乌萨奇能够通过环境微弱变化判断潜在资源点，
并以极短反应时间完成路径选择。
针对草丛中可能出现的未知风险，
采用“先动起来、边跑边判断”的非线性策略，
有效提升了小队在复杂开放环境下的探索效率。
该项目在多次劳动任务中表现稳定，
尤其适用于奖励未知、风险未知但零食可能存在的高不确定性场景。

% ---------------------------------------------------------
% 项目 2
% ---------------------------------------------------------

\subsection{
    三人小队协同探索与情绪稳定支持系统
    \hfill
    小队协作项目
    \hfill
    2025年6月-2026年6月
}

\hspace*{1.5em}
围绕吉伊卡哇、小八与乌萨奇三人小队的日常协同需求，
乌萨奇承担了气氛激活、风险试探与行动破局等关键角色。
在队友犹豫、紧张或陷入低情绪状态时，
能够通过高能量动作、短促口号和不可预测行为打破停滞局面，
为小队提供额外行动动力。
在复杂任务中，
乌萨奇常以非传统决策方式率先接触未知对象，
虽然路径难以复现，
但多次在关键节点触发有效进展。
该项目显著增强了小队面对突发事件时的韧性与连续行动能力。
乌拉亚哈亚哈乌拉乌拉亚哈亚哈乌拉乌拉亚哈亚哈乌拉乌拉亚哈亚哈乌拉乌拉亚哈亚哈乌拉乌拉亚哈亚哈乌拉。

% ---------------------------------------------------------
% 项目 3
% ---------------------------------------------------------

\subsection{
    零食资源发现、评估与即时转化工程
    \hfill
    快乐经济项目
    \hfill
    2025年6月-2026年6月
}

\hspace*{1.5em}
针对街区、洞穴、商店与未知空间中随机出现的食物资源，
乌萨奇形成了一套高度个人化的快速识别与即时处理流程。
该流程包括外观判断、气味确认、风险忽略、
情绪放大与即时消费等步骤，
能够将普通零食迅速转化为高强度快乐输出。
在资源评估阶段，
乌萨奇不依赖复杂表格与长期规划，
而是通过经验、直觉和兴奋程度进行综合判断。
该项目有效提升了个体幸福感，
并在部分情况下带动队友共同进入积极状态，
是吉伊卡哇世界中极具代表性的快乐生产机制。

% ---------------------------------------------------------
% 项目 4
% ---------------------------------------------------------

\subsection{
    未知洞穴与异常事件快速响应机制
    \hfill
    野外探索项目
    \hfill
    2025年6月-2026年6月
}

\hspace*{1.5em}
针对洞穴、森林及未知建筑中频繁出现的异常事件，
乌萨奇参与建立了以高机动侦察、
快速试错和即时撤退为核心的响应机制。
在探索过程中，
乌萨奇能够在缺乏完整情报的前提下迅速进入现场，
通过跳跃、冲刺和近距离观察获取第一手信息。
面对难以解释的道具、生物或场景变化时，
通常采取直接互动策略，
以最短路径验证风险边界。
尽管该机制具有一定不可控性，
但在多次事件中成功推进了小队探索进程，
并体现出极强的环境适应能力。

% ---------------------------------------------------------
% 项目 5
% ---------------------------------------------------------

\subsection{
    零食获取与快乐传播机制研究
    \hfill
    日常行为观察项目
    \hfill
    2025年6月-2026年6月
}

\hspace*{1.5em}
围绕吉伊卡哇世界中零食、奖励与情绪反馈之间的关系，
乌萨奇长期开展非正式行为实验。
在发现可食用资源或疑似奖励物品后，
能够迅速完成接近、确认、展示与消费，
并通过高度夸张的动作和声音将个体快乐扩散至周围环境。
该机制不依赖复杂语言说明，
而是通过表情、动作和短语形成直接感染力。
项目体现出乌萨奇在情绪调动、
现场活跃和小队氛围维护方面的突出能力。

% =========================================================
% 技能特长
% =========================================================

\section{
    \makebox[\widthof{\faWrench}][c]{
        \color{CHU_Blue}{\faWrench}
    }
    \quad
    技能特长
}

\begin{itemize}

    \item
    行动能力：
    高速移动、跳跃突进、危险场景快速试探、复杂地形适应。

    \item
    劳动技能：
    草丛采集、讨伐协助、资源发现、任务现场气氛激活。

    \item
    沟通能力：
    熟练掌握“YAHA”“URA”“哈？”等高密度短语表达，
    具备强辨识度非语言沟通能力。

    \item
    协作能力：
    长期与吉伊卡哇、小八共同活动，
    能够在小队探索、劳动与危机处理中承担破局角色。

\end{itemize}

% =========================================================
% 其他经历
% =========================================================

\section{
    \makebox[\widthof{\faStar}][c]{
        \color{CHU_Blue}{\faStar}
    }
    \quad
    其他经历
}

\begin{itemize}

    \item
    长期活跃于吉伊卡哇世界街区、草丛、洞穴与各类临时任务现场，
    具备跨场景生存与即时适应能力。

    \item
    多次参与小队日常冒险，
    在未知食物、异常道具和突发状况面前保持高度兴奋与持续行动。

    \item
    具备强烈个人风格，
    表达简短但存在感极高，
    适合担任小队气氛发动机、风险试探员与快乐传播节点。

\end{itemize}

\end{document}